\section{DeepSeekMath 与 GRPO}

\subsection{Data Scaling：数学数据不等于 arXiv 论文}

DeepSeekMath 的第一项贡献是数据工程：从大规模网页中识别高质量数学相关内容，而不是简单扩大 arXiv 文本。课程强调，形式化公式密集不一定等于有助于解题；包含自然语言推导、例题和解法的网站可能更有训练价值。

\videofigure{figures/deepseekmath.jpg}{DeepSeekMath-7B 通过数据与 RL 获得强数学表现。}{00:41:10--00:42:06}

\videofigure{figures/math-data.jpg}{高质量数学网页数据与 practical RL 构成方法主线。}{00:42:06--00:43:34}

\videofigure{figures/data-results.jpg}{数据来源选择会显著影响数学推理，盲目加入 arXiv 甚至可能伤害性能。}{00:43:34--00:44:24}

\begin{importantbox}{Data scaling 的单位应是有效学习信号}
更多 token 不等于更多能力。数据是否包含问题--推导--答案结构、难度覆盖、错误纠正与多样表达，通常比“领域关键词比例”更重要。
\end{importantbox}

\subsection{PPO 的 critic 成本}

标准 PPO 通常维护 policy、reference policy、reward model 和 value/critic model，多份大模型占用显存并增加同步成本。

\videofigure{figures/ppo-memory.jpg}{大模型 PPO 的 policy、reference、reward 与 critic 带来高内存压力。}{00:44:24--00:45:12}

\subsection{GRPO：用组内相对奖励替代 critic}

Group Relative Policy Optimization 对同一问题采样一组 $G$ 个回答，用组内均值与标准差构造 advantage：

$$
A_i=\frac{r_i-\mu_G}{\sigma_G+\epsilon},
\qquad
\mu_G=\frac{1}{G}\sum_{j=1}^{G}r_j.
$$

\videofigure{figures/grpo.jpg}{GRPO 用同题候选组作为 baseline，不再训练独立 critic。}{00:45:12--00:46:02}

直觉是：reward 的绝对尺度可能因题目难度变化，但同一道题内“谁比同组更好”更稳定。正确候选得到正 advantage，错误候选得到负 advantage。

GRPO 仍使用 clipped policy update 和 reference KL：

$$
J(\theta)=\mathbb E_i\left[
\min(\rho_i A_i,\operatorname{clip}(\rho_i,1-\epsilon,1+\epsilon)A_i)
-\beta D_{\mathrm{KL}}(\pi_\theta\|\pi_{\mathrm{ref}})
\right].
$$

\videofigure{figures/grpo-gradient.jpg}{多种 RL 方法都可理解为对 $\nabla\log p(\text{output}\mid\text{question})$ 乘以不同梯度系数。}{00:46:02--00:52:46}

\subsection{Online RL 优于固定离线样本}

如果训练始终使用旧 policy 生成的数据，随着 policy 更新，样本分布逐渐过时。Online RL 持续从当前模型采样，reward 和 gradient 更贴近当前失败模式。

\videofigure{figures/online-rl.jpg}{从当前 policy 在线采样优于反复使用旧模型生成的离线轨迹。}{00:52:46--00:53:34}

\begin{warningbox}{组内 advantage 可能为零}
若同一组所有回答都正确或都错误，组内 reward 方差为零，没有有效排序信号。随着模型变强，容易题会频繁出现全对组；难题则可能全错。DAPO 的 dynamic sampling 正是为解决这一浪费。
\end{warningbox}

\subsection{本章小结}

DeepSeekMath 表明 train-time scaling 依赖数据与 RL 系统一起优化。GRPO 用组内相对 reward 移除 critic，降低显存与工程复杂度，但长 reasoning 训练仍会出现熵坍塌、无效 batch 和长度偏差。

